\chapter{The Fog Road}

\section*{Told by}
Mist-walker Kell, an old Ubral woman who lives in a cairn and never lights a fire. She says the fog is not a thing. It is a question. She says fire makes the question impatient. The patient question is the one you can answer. The impatient question eats.

\section*{The Tale}

A child got lost on the moor. The fog came down thick and gray—not rolling, not drifting, but lowering like a lid. She could not see her feet. She could not see her hands when she raised them to her face. She walked for hours, but always returned to the same stone. The stone was flat and warm, though there was no sun. She sat on it. She waited.

A voice spoke from the mist. Not loud. Not close. It came from everywhere at once, like the sound of a bell that had been struck a long time ago and was still ringing.

\textit{``Who walks on my road?''}

"I do," said the child. Her voice was small. The fog did not swallow it. The fog carried it, further than she expected, as if the moor were listening.

\textit{``Are you lost?''}

"I am lost," she said.

\textit{``Are you?''} said the voice. \textit{``You are walking exactly where I need you. You have not missed any turns. You have not taken any wrong roads. You have simply not arrived yet. That is not the same as being lost. That is called traveling. But you still have not answered my first question. Who walks?''}

The child thought. She could not give her name—names are for daytime, for hearths, for people who know where they belong. She could not say her clan—clans are for fires, and there were no fires here. She could not say where she was from, because the fog had erased the direction.

She thought of what her grandmother had told her, once, when the fog came down around their wagon: \textit{``The moor does not take those who belong to it. It only takes those who argue.''}

"Someone who belongs to the fog," she said.

The voice laughed. Not a cruel laugh. A surprised laugh. A laugh that said \textit{you are the first person to say that in a hundred years}.

\textit{``Then you may stay. Not forever. Just until you find the stones.''}

The fog parted. Not all at once—slowly, like a curtain drawn by a patient hand. The child saw a path of white stones, each one placed just far enough from the next that she had to stretch. She followed it home.

At the edge of the moor, she looked back. The fog was gone. The moor was ordinary grass and heather. She did not tell anyone what she had seen. She did not need to. Her grandmother knew.

"You walked with the Urisk," her grandmother said. "It asks every traveler the same question. Most answer with their names. Then it keeps them. It does not take their lives. It takes their \textit{arrival}. They walk and walk and never reach the edge. They are still walking. They have been walking for centuries. They are not unhappy. They are just... not home."

The child touched the stone in her pocket. She did not remember picking it up. It was warm. It was flat. It was the same stone.

She kept it for seventy years. When she died, her daughter placed it back on the moor. The fog came down that night. It has not lifted since.

\section*{The Witch's Closing}
\textit{``The fog does not hide. It asks. It asks if you belong. If you do, it will guide you. If you do not, it will keep you—not cruelly, not kindly, just... permanently. The moor has a long memory. It remembers every answer it has ever heard. It is still waiting for the right one.''}

\section*{What Lurks in the Shadow of the Tale}

\subsection*{The Urisk of the Moor}

A spirit of the fog. It appears to lost travelers and asks, \textit{``Who walks on my road?''} It does not repeat the question. It does not clarify. It asks once, and the answer decides everything. If the traveler gives a name—any name, true or false—the Urisk learns it and may use it to compel the traveler to stay. Not to hurt them. To \textit{keep} them. The moor is lonely. The Urisk is old. Loneliness is not malice, but it is not kindness either.

If the traveler claims belonging to the fog—not ownership, not mastery, but belonging, the way a stone belongs to the hill, the way rain belongs to the sky—the Urisk guides them safely and may become an ally. It will not appear when called. It will not answer questions. But the fog will never hide the path from that traveler again.

\paragraph{Tags / Mechanics:}
\begin{itemize}
\item [SPIRIT] [THRESHOLD] [TRUTH]
\item To negotiate with the Urisk, a character must succeed on a Spirit + Lore test (DV 4) while offering a truth about themselves that is not their name. A memory. A regret. A hope. Something the fog can hold. On success, the Urisk parts the fog, and the character may re-roll any failed navigation check on the moor for the rest of the journey. On failure, the Urisk does not harm them. It simply says, \textit{``Not yet. Walk further.''} And the fog does not lift.
\item If the character gives their name—even in panic, even accidentally—the Urisk can later appear at a dramatic moment, offering a choice: a path that seems easy but leads to danger, a shortcut that costs more than it saves, a door that opens onto a room the character has already left. The GM gains 3 SB to spend on complications related to \textit{getting lost}. The character does not need to accept the Urisk's offered path. But they will see it. They will always see it.
\item A character who successfully claims belonging to the fog gains the following: once per session, when they are lost (geographically, morally, or narratively), they may ask the GM, \textit{``Which way feels like home?''} The GM answers honestly, though not always clearly. The fog remembers its own.
\end{itemize}

\paragraph{Adventure Hook:}
A shepherd's child has vanished into the fog. The villagers are afraid to search—not because they are cowards, but because three other children vanished in the same fog twenty years ago, and none of them came back. The party must find the Urisk and answer its question. But the fog is thick, and the moor is wide, and the Urisk has heard a great many names. It is still waiting for the right one. It is not impatient. It has never been impatient. That is what makes it dangerous.

\section*{Colophon}
Cairn on the high moor. The old woman did not light a candle. She did not need one. The fog was bright enough.

\textit{``Light attracts the wrong things,''} she said. \textit{``The dark is easier to negotiate with. In the dark, you can hear the questions before you see the teeth. In the light, the teeth are all you see.''}

She poured a cup of water onto the stone floor. The water did not pool. It vanished.

\textit{``The moor is thirsty,''} she said. \textit{``It has been thirsty for a long time. The fog is just its breath. Do not hold your breath back. The moor does not like to wait.''}